\section{Discussion and conclusion}

The standard honest/adversarial partition is sufficient only when “honest” also means “independently correct for the fault being exploited.” Correlated implementation failure separates those concepts.

The resulting JAM/ELVES analysis has three nested boundaries.

First, the dispute layer has an attribution boundary. If the combined adversarial-plus-buggy-positive bloc exceeds roughly one third of the validator population, independently correct validators no longer suffice by themselves to construct a bad verdict. The report can still be stopped through a wonky outcome, but the ordinary Gray Paper offender-attribution path changes. This is a protocol-rule result, not a branching-model result.

Second, the audit layer has a correction boundary. Under the ELVES pessimistic branching extension,
\[
\lambda_x=Fh_x,
\]
and correction ceases to be supercritical at \(h_x=1/F\). For \(F=2\), the corresponding wrong-positive boundary is one half, above the one-third attribution threshold.

Third, sampled correction introduces dilution. With
\[
\lambda_x=F\frac{C_x}{N},
\]
validator growth can reduce corrective reproduction when new validators do not increase the independently correct set for the relevant fault. This can occur even if every entrant is honest and the adversarial fraction falls.

Together these results suggest that a useful security variable is not raw validator count or raw client count but fault-specific independent corrective capacity. Because future faults are unknown, operators cannot observe that capacity directly. They can only monitor exposure to plausible common components and preserve diversity margins across multiple layers.

The most concise summary is therefore:
\[
\boxed{
\text{honest stake}
\neq
\text{independently corrective stake}
\neq
\text{attacker-attributable stake}.
}
\]

The first inequality matters for auditing. The second matters for economics. The Gray Paper and ELVES already provide a strong framework for reasoning about adversarial and availability failures; correlated honest failure asks where shared implementation dependencies enter that framework.

The immediate next step is not another model extension. It is review by the ELVES and JAM authors, especially on the wonky-verdict economics and the substitution of \(h_x\) into the branching proof. If those mappings are accepted, fault-domain exposure becomes a quantitatively defined protocol margin rather than only a qualitative argument for client diversity. If they are rejected, the reasons will identify exactly which additional mechanism closes the gap.
